\section{Sequence properties, restriction digests and file formats}
This section checks four bio/ libraries against Biopython 1.88 ProtParam + Restriction and gffutils. Three pass exactly. One, the restriction-site finder, had a real bug that we fixed.
\subsection{Protein physicochemical properties}
On 500 UniProtKB reviewed human proteins (UniProtKB reviewed human, length 50-2000, first 500 of search (rest.uniprot.org, 2026-09-25)), every property matches ProtParam (Table~\ref{tab:pp}).
\begin{table}[h]\centering\small\begin{tabular}{lrrrr}\hline property & $n$ & within tol. & max abs diff & Pearson\\\hline
mw & 500 & 500 (rel 1e-06) & 0 & 1.0000000000\\
gravy & 500 & 500 (rel 1e-06) & 0 & 1.0000000000\\
aromaticity & 500 & 500 (rel 1e-06) & 2.8e-17 & 1.0000000000\\
instability & 500 & 500 (rel 1e-06) & 0 & 1.0000000000\\
pI & 500 & 500 (abs 0.01) & 5.9e-05 & 0.9999999999\\
ext\_reduced & 500 & 500 (rel 1e-06) & 0 & 1.0000000000\\
ext\_oxidized & 500 & 500 (rel 1e-06) & 0 & 1.0000000000\\
\hline\end{tabular}\caption{bio.proteinprops vs Biopython ProtParam (benchmarks/sweep\_proteinprops\_restriction.json).}\label{tab:pp}\end{table}
\paragraph{(F36) Average molecular weight.} $M=\sum_a n_a m_a-(L-1)\,m_{\mathrm{H_2O}}$, with $m_a$ the average residue mass (free amino acid) and one water lost per peptide bond.
\paragraph{(F37) Net charge.} $Z(\mathrm{pH})=\sum_{i\in +}\frac{n_i}{1+10^{\mathrm{pH}-pK_i}}-\sum_{j\in -}\frac{n_j}{1+10^{pK_j-\mathrm{pH}}}$ (Henderson--Hasselbalch per ionisable group, termini included).
\paragraph{(F38) Isoelectric point.} $\mathrm{pI}$ is the root of $Z(\mathrm{pH})=0$; $Z$ is strictly decreasing in pH, so bisection on $[0,14]$ converges, and the residual $6\times10^{-5}$ difference is the bisection tolerance.
\paragraph{(F39) Instability index.} $II=\frac{10}{L}\sum_{i=1}^{L-1}\mathrm{DIWV}(x_i,x_{i+1})$ (Guruprasad 1990), and GRAVY $=\frac1L\sum_i h_{KD}(x_i)$ (Kyte--Doolittle).
\paragraph{(F40) Extinction coefficient at 280 nm.} $\varepsilon=5500\,n_W+1490\,n_Y+125\,n_{\mathrm{cystine}}$ M$^{-1}$cm$^{-1}$, with $n_{\mathrm{cystine}}=\lfloor n_C/2\rfloor$ in the oxidised form and $0$ when reduced.
\subsection{Restriction sites: a bug in degenerate recognition sequences}
We digested pBR322 (J01749.1, circular, 4,361 bp) and phage lambda (NC\_001416.1, linear, 48,502 bp) with all 610 enzymes and compared cut positions with Biopython Restriction. Before the fix (sugarcode-ai 699ffee), 533 and 529 enzymes agreed. Two causes explained the disagreements. First, the reverse-strand pattern for a non-palindromic-looking site was built with a plain DNA reverse complement that leaves IUPAC codes unchanged. AccI (GT\^{}MKAC) is in fact palindromic under IUPAC complementation (M$\leftrightarrow$K), but the plain reverse complement GTKMAC was searched as a second pattern, which matches GTGAAC and GTTCAC, neither an AccI site. On pBR322 this gave 5 AccI sites where Biopython finds 2. Second, on linear DNA the finder reported Type IIS cuts that fall outside the molecule. After the fix (sugarcode-ai fe2185d), 608 and 608 enzymes agree. The remaining two, LpnPI, MspJI, are modification-dependent enzymes that do not cut unmethylated DNA; we leave them open.
\paragraph{(F41) IUPAC complement.} The map $c$: A$\leftrightarrow$T, C$\leftrightarrow$G, R$\leftrightarrow$Y, K$\leftrightarrow$M, B$\leftrightarrow$V, D$\leftrightarrow$H, with S, W, N fixed, is the unique involution on the 15 codes that complements every base in the set a code denotes. A site $s$ is palindromic iff $s=\overline{c(s)}$ (reversed), which for AccI holds only under $c$, not under the four-letter map.
\subsection{Stockholm alignments}
bio.stockholm parsed 15 Rfam seed alignments (4,471 sequences). 15/15 match Biopython on sequence identifiers and order, aligned sequences and SS\_cons, and survive a write--parse round trip unchanged.
\begin{table}[h]\centering\small\begin{tabular}{llrrc}\hline file & AC & seqs & width & all checks\\\hline
RF00001.sto & RF00001 & 712 & 230 & yes\\
RF00005.sto & RF00005 & 954 & 118 & yes\\
RF00010.sto & RF00010 & 458 & 996 & yes\\
RF00004.sto & RF00004 & 208 & 278 & yes\\
RF00015.sto & RF00015 & 177 & 261 & yes\\
RF00017.sto & RF00017 & 91 & 412 & yes\\
RF00023.sto & RF00023 & 477 & 954 & yes\\
RF00050.sto & RF00050 & 146 & 221 & yes\\
RF00059.sto & RF00059 & 115 & 211 & yes\\
RF00162.sto & RF00162 & 457 & 239 & yes\\
RF00167.sto & RF00167 & 133 & 113 & yes\\
RF00168.sto & RF00168 & 47 & 274 & yes\\
RF00174.sto & RF00174 & 434 & 461 & yes\\
RF00234.sto & RF00234 & 18 & 236 & yes\\
RF00504.sto & RF00504 & 44 & 215 & yes\\
\hline\end{tabular}\caption{bio.stockholm vs Biopython AlignIO (benchmarks/sweep\_stockholm\_biopython.json).}\end{table}
\paragraph{(F42) Pairwise identity.} $\mathrm{PID}(x,y)=\#\{i: x_i=y_i,\ x_i,y_i\notin\{-,.\}\}\,/\,\#\{i: x_i,y_i\notin\{-,.\}\}$, case-insensitive; our implementation equals an independent count on every file.
\subsection{GFF3 annotation}
On the RefSeq E. coli K-12 annotation (9,523 records), bio.gff and gffutils 0.14 (sqlite, in-memory) agree on every record (9,523/9,523) and on the children of 300/300 genes, and write--parse is lossless. A first comparison showed two child mismatches; both were CDS features split over two lines with a shared ID, which gffutils renames X\_1. That was a comparison artifact, not a sugarcode error, and we record it so the numbers are traceable.
\begin{table}[h]\centering\small\begin{tabular}{lr}\hline feature type & records\\\hline
gene & 4506\\
CDS & 4340\\
exon & 216\\
pseudogene & 145\\
ncRNA & 108\\
tRNA & 86\\
mobile\_genetic\_element & 50\\
sequence\_feature & 48\\
\hline\end{tabular}\caption{Feature types in GCF\_000005845.2 (benchmarks/sweep\_gff\_gffutils.json).}\end{table}
\subsection{Newick trees}
bio.newick parsed 6 Open Tree of Life synthetic subtrees with 3,655 leaves in total, including quoted labels with spaces and parentheses and empty leaves left by the height limit. Leaf sets equal Biopython Bio.Phylo and DendroPy, internal-node counts and internal labels equal Bio.Phylo, and write--parse is lossless in 6/6 trees. These trees have no branch lengths, so this checks topology and labels only.
\begin{table}[h]\centering\small\begin{tabular}{lrrrc}\hline root & leaves & unnamed & internal & all checks\\\hline
Insecta\_ott1062253 & 1094 & 14 & 248 & yes\\
Laurasiatheria\_ott392223 & 728 & 13 & 118 & yes\\
Carnivora\_ott44565 & 53 & 13 & 48 & yes\\
Metazoa\_ott691846 & 1581 & 19 & 179 & yes\\
Aves\_ott81461 & 142 & 16 & 80 & yes\\
Primates\_ott913935 & 57 & 11 & 57 & yes\\
\hline\end{tabular}\caption{bio.newick vs Bio.Phylo/DendroPy (benchmarks/sweep\_newick\_biophylo.json).}\end{table}
\paragraph{(F43) Tree size identity.} For a rooted tree with $n$ leaves, $m$ internal nodes and out-degrees $d_v$, the edge count is $n+m-1=\sum_{v\,\mathrm{internal}}d_v$; a strictly binary tree has $m=n-1$. The Carnivora subtree has 53 leaves and 48 internal nodes, so it is far from binary, which exercises the parser on polytomies.
\subsection{SAM records and a second bug}
bio.sam parsed 10,943 real reads (1000 Genomes phase 3 NA12878 chrom20 low-coverage BAM, 20:1000000-1200000). Name, flag, reference, position, mapping quality, CIGAR operations, mate fields, sequence, qualities and typed tags equal pysam for all 10,943 reads. The derived reference end did not: 10,905/10,943 before the fix. The 38 failures were unmapped mates (flag 0x4) that bwa places at the mapped mate's position and gives a CIGAR. The SAM specification says those fields are not an alignment, and htslib reports no end; sugarcode computed one. After sugarcode-ai 87cf78e all 10,943 agree.
\paragraph{(F44) Reference span of a CIGAR.} $\mathrm{end}=\mathrm{POS}+\sum_{(n,o)\in C,\ o\in\{M,D,N,=,X\}}n-1$ for mapped reads (1-based, inclusive); query length is $\sum_{o\in\{M,I,S,=,X\}}n$. For flag $\&\,4\ne0$ the end is undefined.
\subsection{FASTQ reads}
On the first 20,000 reads of ENA run SRR622461 (the same NA12878 library), bio.fastq equals Biopython on identifiers (20,000), sequences (20,000) and per-base Phred scores (20,000), detects offset 33, and round-trips losslessly. GC was 0.4252 against 0.4252 (Biopython, ambiguous bases removed) and 0.422 (ambiguous bases kept in the denominator). The difference is a convention: sugarcode divides by A+C+G+T. We document it rather than change it.
\paragraph{(F45) Phred score.} $Q=-10\log_{10}P_{err}$, encoded as ASCII $Q+33$; the mean read quality reported is the arithmetic mean of $Q$, not $-10\log_{10}$ of the mean error probability, which would be lower for reads with a few bad bases.
\subsection{PDB and mmCIF structures: a third bug}
We parsed 30 RCSB entries in both PDB and mmCIF format (94,121 atom sites, all alternate locations and models) and compared every atom with gemmi 0.7.5. The PDB-format parser was identical in 30/30 entries. The mmCIF parser failed on all 30 real files (error: loop\_ row count 3 is not a multiple of 4 columns (all 30 files)). Two causes: POSIX shell-style splitting broke atom names containing a prime (O5$'$, C1$'$), and loops were closed early when a row contained an unquoted item name as a value, which RCSB files do in the revision-history tables. A CIF 1.1 tokenizer (a quote closes only when followed by whitespace) and closing loops only at row boundaries fixed both (sugarcode-ai fc1fd1e); mmCIF is now identical to gemmi in 30/30 entries. Unit tests with toy CIF strings had passed throughout. Only real files exposed the bug.
\paragraph{(F46) Loop row invariant.} A loop with $c$ columns and $t$ value tokens is well formed iff $t\equiv0\pmod c$; closing the loop at the first reserved-looking token is only safe when the running count $t\bmod c=0$, which is the rule we now apply.
\subsection{Pileup}
Using the same reads, bio.pileup and pysam 0.24.1 pileup (stepper nofilter, flag\_filter=4, no quality filters, overlaps kept) agree on base counts at 197,437/197,437 positions and on deletion counts at 197,437/197,437. sugarcode also writes deletions as `*' in the base counts, following samtools mpileup; the comparison maps these to the deletion count.
\paragraph{(F49) Column depth.} $D(p)=\sum_{r}\mathbb 1[p\in[\mathrm{POS}_r,\mathrm{end}_r]\ \wedge\ \mathrm{op}_r(p)\in\{M,=,X,D\}]$, split into base-covering reads ($M,=,X$) and deletions ($D$); reference skips ($N$) and clipped bases do not count.
\subsection{Transcription-factor motifs}
We read 10 JASPAR 2024 profiles and scanned BRCA1 (10 kb) and phage lambda (48.5 kb) on both strands, comparing with Biopython 1.88 Bio.motifs (jaspar reader, pseudocount 0.5, uniform background, log2). Log-odds matrices agree to $2\times10^{-15}$, per-window scores to $4\times10^{-6}$ (Biopython stores scores in single precision), and the hit sets at 80\% of the score range are identical for every motif and sequence.
\begin{table}[h]\centering\small\begin{tabular}{lrrr}\hline motif & length & hits BRCA1 & hits lambda\\\hline
MA0003.4 TFAP2A & 14 & 27 & 61\\
MA0079.3 SP1 & 11 & 235 & 599\\
MA0080.4 SPI1 & 14 & 16 & 61\\
MA0099.3 FOS::JUN & 10 & 77 & 286\\
MA0105.4 NFKB1 & 13 & 29 & 52\\
MA0106.3 TP53 & 18 & 8 & 21\\
MA0108.2 TBP & 15 & 154 & 530\\
MA0112.3 ESR1 & 17 & 11 & 18\\
MA0113.3 NR3C1 & 17 & 2 & 36\\
MA0139.1 CTCF & 19 & 13 & 40\\
\hline\end{tabular}\caption{Motif hits at 80\% of range, identical to Biopython (benchmarks/sweep\_motif\_biopython.json).}\end{table}
\paragraph{(F50) Log-odds score and relative threshold.} $S(w)=\sum_{i=1}^{L}\log_2\frac{(c_{i,w_i}+\alpha)/(N_i+4\alpha)}{q_{w_i}}$ with pseudocount $\alpha=0.5$ and background $q=1/4$; a window is a hit when $S(w)\ge S_{\min}+f\,(S_{\max}-S_{\min})$, $f=0.8$, where $S_{\max}=\sum_i\max_b$ and $S_{\min}=\sum_i\min_b$ of the log-odds matrix.
\subsection{BED intervals: a fourth bug}
We fetched two hg38 tracks from the UCSC REST API: CpG islands on chr20 (847 intervals) and RepeatMasker on chr20:0--5\,Mb (11,699 intervals), wrote them as BED and compared bio.bed with bioframe 0.8.0 (merge, overlap). Parsing and 600 random overlap queries agreed exactly (300/300 and 300/300). merge\_intervals did not: its docstring promised bedtools-style merging, yet it refused to join book-ended intervals. RepeatMasker has 3,112 book-ended pairs, so sugarcode returned 11,364 merged intervals where bedtools/bioframe return 8,251. CpG islands have no book-ended pairs and hid the bug (847 = 847). We added a bedtools-compatible min\_dist argument (default 0; $-1$ for overlap-only) and a matching CLI flag; after the fix both tracks match bioframe exactly (8,251 = 8,251).
\paragraph{(F51) Merge rule.} For intervals sorted by start with running end $e^*$, interval $[s,e)$ joins the current block iff $s\le e^*+d$; $d=0$ joins book-ended intervals ($s=e^*$) and $d=-1$ requires at least one shared base ($s<e^*$).
\subsection{GenBank flat files: a fifth bug}
We fetched 25 GenBank records one by one from NCBI nuccore (phage, viral, mitochondrial, bacterial, RefSeqGene and 13 human RefSeq mRNAs) and compared bio.genbank with Biopython 1.88 SeqIO genbank on 1,867 features. Sequences and feature counts already agreed, but only 1,829/1,867 feature locations and 5/302 CDS /translation qualifiers matched. The parser appended every continuation line of a multi-line qualifier to the feature location instead of to the qualifier: /translation kept its first line only, and when a continuation began with digits the location itself was corrupted (E.~coli thrA CDS: 337..2799 became 337..27991). Partial locations (<1..>200) and single bases produced no spans at all. We rewrote qualifier continuation (INSDC joining: no spaces for /translation, one space otherwise, doubled quotes unescaped), added a location parser for join/order, complement, partial ends, between-base sites and remote parts, and kept all IUPAC letters in ORIGIN so coordinates cannot shift. After the fix: spans 1,867/1,867, strands 1,867/1,867, /translation 302/302, CDS extraction 302/302.
The fix exposed a downstream case: KCNQ1's record NG\_016178.2 carries only a partial CDS, complement(<42918..>43038). The old parser read no spans and bio.splice reported 'no CDS spans'; with spans read, the junction map came back 'ok' but empty. We made it fail loud with status 'partial CDS in record'. The splice-region discovery pipeline reads hg38 directly and does not use this parser, so its results are unchanged.
\paragraph{(F52) Qualifier continuation.} For a qualifier whose value spans lines $\ell_1,\dots,\ell_m$, $v=\ell_1\oplus_\sigma\ell_2\oplus_\sigma\cdots\oplus_\sigma\ell_m$ with separator $\sigma=\varnothing$ for /translation and $\sigma=$ one space otherwise; a line starting with / opens a new qualifier only when the running count of quote characters in $v$ is even.
\subsection{Core sequence operations}
On the same 25 GenBank records, bio.sequence.translate reproduced the /translation qualifier for 285/285 standard-code CDSs (17 CDSs with other genetic codes, codon\_start$\ne$1, transl\_except or ribosomal slippage were set aside) and matched Bio.Seq.translate on all 285. Reverse complements agreed on 25/25 whole records, IUPAC motif search on 192/192 (motif, record) pairs, Wallace $T_m$ on 500/500 lambda oligos (8--13 nt), and ORF calls ($\ge$100 aa, both strands) on 22/22 records (239 = 239 ORFs) against an independent Biopython frame scan. GC fraction agreed to $10^{-12}$ on 24/25; the exception is the human mitochondrial genome, whose single N sugarcode leaves out of the denominator. Protein masses on 500 UniProt proteins differ from ProteinAnalysis by at most 5.0e-05 (relative), from two-decimal residue masses.
Reading the code for this comparison turned up a latent bug the data could not: clean\_dna silently dropped IUPAC ambiguity letters, so ATGRCTTAA translated as ML (frame shifted) instead of MX. No CDS in this set carries an ambiguity code, so the real-data numbers are identical before and after; ambiguity letters now become N (sugarcode 143d1c4), with a regression test. We report it as a code-review finding, not a data finding.
\paragraph{(F53) Frame preservation.} Cleaning must be a letter-wise map $c:\Sigma\to\{A,C,G,T,N\}\cup\{\epsilon\}$ with $c(x)=\epsilon$ only for non-letters; then codon $k$ of the cleaned sequence is $c(s_{3k})c(s_{3k+1})c(s_{3k+2})$ and every downstream codon keeps its frame.
